\documentclass[a4paper,10pt]{scrartcl}

\usepackage[ngerman]{babel}
\usepackage[utf8]{inputenc}
\usepackage[T1]{fontenc}
\usepackage[babel,german=quotes]{csquotes}
\usepackage{hyperref}
\usepackage{geometry}
\usepackage{graphicx}
\usepackage{enumitem}
\usepackage{multicol}

\setlength\parindent{0pt}
\setitemize{itemsep=0pt}


\geometry{
	a4paper,
	left=20mm,
	top=20mm,
	bottom=15mm,
	footskip=5mm
}

\subject{Gruppe \textit{Nummer}}
\title{Pflichtenheft für das Thema \glqq\textit{Titel}\grqq}
\subtitle{Betreuer: \textit{Name}}
\author{\textit{Namen der Teilnehmenden}}
\date{\textit{Datum}}

\begin{document}

\maketitle

\textit{\textbf{Anmerkung:} Alles, was in diesem Dokument kursiv geschrieben ist, muss von der Gruppe entsprechend ihres Themas angepasst (und diese Anmerkung entfernt) werden}

\section{Zielbestimmung}

\subsection{Muss-Kriterien}

\begin{itemize}
    \item \textit{Stichpunktartige Auflistung der verpflichtenden Anforderungen}
\end{itemize}

\subsubsection{Halbzeit-Ziele}
\begin{itemize}
    \item \textit{Stichpunktartige Auflistung der Muss-Kriterien, die zur Hälfte des Praktikums erfüllt sein sollen (muss mindestens 50\% der Muss-Kriterien entsprechen)}
\end{itemize}

\subsubsection{Zuständigkeiten}
\begin{itemize}
    \item \textit{Stichpunktartige Auflistung der groben Zuständigkeiten (anhand der Muss-Kriterien) für jeden Teilnehmenden}
\end{itemize}

\subsection{Kann-Kriterien}

\begin{itemize}
    \item \textit{Stichpunktartige Auflistung der optionalen Anforderungen}
\end{itemize}

\subsection{Abgrenzungskriterien}

\begin{itemize}
    \item \textit{Optionale, stichpunktartige Auflistung von Funktionalitäten, die explizit nicht umgesetzt werden müssen}
\end{itemize}


\section{Einsatz}

\subsection{Anwendungsbereiche}

\textit{Kurze Beschreibung (Stichpunkte oder 1 bis 2 Sätze), wo das Projekt-Ergebnis eingesetzt werden soll (z.B. Hochschulen oder Krankenhäuser)}

\subsection{Zielgruppe}

\textit{Kurze Beschreibung (Stichpunkte oder 1 bis 2 Sätze), von wem das Projekt-Ergebnis verwendet werden soll (z.B. Studierende, Lehrende oder Ärztinnen und Ärzte)}


\section{Umgebung}

\subsection{Software}

\begin{itemize}
    \item \textit{Stichpunktartige Auflistung der verwendeten Programmiersprachen, Frameworks und Tools jeweils mit Versionsnummer und Links zu den zugehörigen Internetseiten}
\end{itemize}

\subsection{Hardware}

Die Hardware muss alle Anforderungen der im vorherigen Punkt genannten Software erfüllen.

\textit{Optional: Kurze Beschreibung (Stichpunkte oder 1 bis 2 Sätze) weiterer Anforderungen, wie z.B. eine erforderliche Internet-Verbindung, oder bestimmte Geräte}


\section{Funktionalität}

\textit{Kurze Beschreibung (Stichpunkte oder 3 bis 4 Sätze) von typischen Anwendungsfällen/Arbeitsabläufen mit dem Projekt-Ergebnis}


\section{Daten}

\textit{Kurze Beschreibung (Stichpunkte oder 3 bis 4 Sätze) von verwendeten und erzeugten Daten}


\section{Leistungen}

Die \textit{kurze Auflistung der geplanten Funktionalitäten} soll weitestgehend performant sein und so in einer angemessenen Zeit geschehen.


\section{Benutzungsoberfläche}

\textit{Dieser Abschnitt kann bei Themen ohne GUI weggelassen werden; ansonsten grobe Skizzen der wichtigsten GUI-Ansichten mit jeweils kurzer Beschreibung (Stichpukte oder 1 bis 2 Sätze pro Skizze)}


\section{Qualitätsziele}

Der Code soll gut kommentiert und klar strukturiert sein, sodass sich neue Entwicklerinnen und Entwickler gut in dem Projekt zurechtfinden und es leicht erweitern und ändern können. Um dies zu erleichtern, wird außerdem ein Developerhandbuch erstellt.\\
Das Ergebnis des Projekts soll eine hohe Benutzerfreundlichkeit haben und für Anwenderinnen und Anwender wird zusätzlich ein Benutzerhandbuch bereitgestellt.

\end{document}
